\documentclass[11pt]{article}

\usepackage{amsmath,amssymb,amsthm,mathtools}
\usepackage[T1]{fontenc}
\usepackage{lmodern}
\usepackage[margin=1in]{geometry}
\usepackage{enumitem}
\usepackage{hyperref}
\usepackage{microtype}

\hypersetup{
  colorlinks=true,
  linkcolor=blue,
  citecolor=blue,
  urlcolor=blue
}

\newtheorem{theorem}{Theorem}[section]
\newtheorem{proposition}[theorem]{Proposition}
\newtheorem{corollary}[theorem]{Corollary}
\newtheorem{definition}[theorem]{Definition}
\newtheorem{lemma}[theorem]{Lemma}
\newtheorem{remark}[theorem]{Remark}

\newcommand{\Z}{\mathbb Z}

\title{Draft graph-theoretic Hamilton decomposition proof for odd $m$}
\author{}
\date{March 2026}

\begin{document}
\maketitle

\begin{abstract}
This note is a working draft of the graph-theoretic proof version for the odd
$m$, $d=5$ torus.

Unlike 099, which packages the odd-$m$ D5 globalization theorem, the present
note aims at the actual graph-level statement:
\begin{quote}
the directed torus $D_5(m)$ admits a decomposition into five arc-disjoint
directed Hamilton cycles.
\end{quote}

The draft is intentionally incomplete. Every missing graph-level bridge,
reduction, or coverage step is left explicitly as \emph{Sorry.} so the document
can be used as an active proof scaffold rather than as an overclaimed theorem
note.
\end{abstract}

\section{Graph-theoretic target}

Let
\[
D_5(m)=\operatorname{Cay}((\Z_m)^5,\{e_0,e_1,e_2,e_3,e_4\})
\]
be the directed $5$-torus.

Its vertex set is
\[
V=(\Z_m)^5,
\]
and its directed arc set is
\[
A=\{(x,x+e_i): x\in V,\ 0\le i\le 4\}.
\]

\begin{definition}[directed Hamilton cycle]
A \emph{directed Hamilton cycle} in $D_5(m)$ is a directed cycle that visits
every vertex of $V$ exactly once.

Equivalently, it is a cyclic permutation of the $m^5$ vertices whose directed
successor arc lies in $A$ at every step.
\end{definition}

\begin{definition}[Hamilton decomposition]
A \emph{Hamilton decomposition} of $D_5(m)$ is a family
\[
H_0,H_1,H_2,H_3,H_4
\]
of five directed Hamilton cycles such that:
\begin{enumerate}[label=\textup{(\arabic*)},leftmargin=2.5em]
\item the arc sets of the $H_i$ are pairwise disjoint;
\item their union is the full directed arc set $A$.
\end{enumerate}

Since $|V|=m^5$ and each Hamilton cycle has exactly $m^5$ arcs, such a
decomposition would contain exactly $5m^5$ arcs in total, which matches the
full directed arc count of $D_5(m)$.
\end{definition}

\begin{theorem}[draft graph-theoretic main theorem]
Let $m$ be odd in the final intended D5 range.

Then the directed torus
\[
D_5(m)=\operatorname{Cay}((\Z_m)^5,\{e_0,e_1,e_2,e_3,e_4\})
\]
admits a Hamilton decomposition into five arc-disjoint directed Hamilton
cycles.
\end{theorem}

\begin{proof}[Proof status]
\textbf{Sorry.} The rest of this note records the missing theorem layers that
still have to be inserted in order for this theorem to become an actual proof
rather than a target statement.
\end{proof}

\section{Imported globalization theorem}

\begin{theorem}[odd-$m$ D5 globalization theorem]
Assume odd $m$ lies in the accepted D5 regime and work in the best-seed channel
\[
R1 \to H_{L1}.
\]
Then every actual lift of the exceptional cutoff row $\delta=3m-3$ continues
through
\[
3m-2 \to 3m-1
\]
and lies in the regular continuing endpoint class there.

Consequently:
\begin{enumerate}[label=\textup{(\arabic*)},leftmargin=2.5em]
\item no mixed-status realized $\delta$ exists;
\item fixed realized $\delta$ determines tail length;
\item fixed realized $\delta$ determines endpoint class and padded future word;
\item $\rho=\rho(\delta)$ globally;
\item raw global $(\beta,\delta)$ is exact on the true accessible boundary
      union.
\end{enumerate}
\end{theorem}

\begin{proof}[Proof status]
Imported from
\texttt{d5\_099\_one\_pass\_odd\_m\_globalization\_package}.
\end{proof}

\section{Reduction from the full odd-$m$ D5 problem}

\begin{theorem}[channel reduction theorem]
To prove Hamilton decomposition of $D_5(m)$ for odd $m$, it is enough to prove
the globalization theorem on the best-seed channel
\[
R1 \to H_{L1},
\]
together with the statement that all other channels are already closed or
reduce to this one.
\end{theorem}

\begin{proof}[Proof status]
\textbf{Sorry.} This theorem needs an explicit channel-classification statement
and a precise proof that the present channel is either the last unresolved case
or a canonical representative of all unresolved cases.
\end{proof}

\begin{remark}[exact missing content]
The missing reduction package should say one of the following, explicitly:
\begin{enumerate}[label=\textup{(\arabic*)},leftmargin=2.5em]
\item every odd-$m$, $d=5$ configuration reduces to the best-seed channel;
\item all other channels are independently solved, and this channel is the
      unique remaining one;
\item the final Hamilton decomposition theorem is a finite disjunction of
      channel theorems, of which the present channel is the only open branch.
\end{enumerate}
At the moment, none of those is written out in this note.
\end{remark}

\section{From raw global exactness to graph-level successor maps}

\begin{definition}[target color factors]
The intended graph-theoretic completion should produce five maps
\[
F_0,F_1,F_2,F_3,F_4 : V \to V
\]
with the following properties:
\begin{enumerate}[label=\textup{(\arabic*)},leftmargin=2.5em]
\item for every $x\in V$ and every $i$, the arc $(x,F_i(x))$ is one of the
      directed torus arcs in $A$;
\item each $F_i$ is a permutation of $V$;
\item the arc sets
      \[
      A_i=\{(x,F_i(x)):x\in V\}
      \]
      are pairwise disjoint and together cover $A$.
\end{enumerate}
When those hold, each $A_i$ is a directed spanning $1$-factor.
\end{definition}

\begin{proof}[Definition status]
\textbf{Sorry.} The exact formula or witness rule that defines the $F_i$ is not
restated here yet. This is only the target shape of the graph-level objects.
\end{proof}

\begin{theorem}[global bridge descends to full-torus successor maps]
Assume:
\begin{enumerate}[label=\textup{(\arabic*)},leftmargin=2.5em]
\item the channel reduction theorem;
\item the odd-$m$ D5 globalization theorem;
\item the accepted D5 witness package that associates boundary / chain data to
      directed torus arcs.
\end{enumerate}
Then the local D5 witness rule descends to globally well-defined maps
\[
F_0,\dots,F_4
\]
on the full torus as in the previous definition.
\end{theorem}

\begin{proof}[Proof status]
\textbf{Sorry.} This is the main missing bridge from the boundary-globalization
package to the graph-theoretic statement. The theorem should explain:
\begin{enumerate}[label=\textup{(\arabic*)},leftmargin=2.5em]
\item how the symbolic/boundary system defines outgoing torus arcs;
\item why raw global exactness removes local ambiguity and makes that choice
      globally well defined;
\item why the resulting maps really live on all vertices of the full torus
      rather than only on the accessible boundary union.
\end{enumerate}
\end{proof}

\begin{proposition}[each $F_i$ uses only torus arcs]
For each $i$ and each $x\in V$, one has
\[
F_i(x)-x \in \{e_0,e_1,e_2,e_3,e_4\}.
\]
\end{proposition}

\begin{proof}[Proof status]
\textbf{Sorry.} This should be the direct graph-level translation of the
accepted D5 local witness / color rule.
\end{proof}

\begin{proposition}[each $F_i$ is a permutation]
For each $i$, the map $F_i:V\to V$ is a bijection.
\end{proposition}

\begin{proof}[Proof status]
\textbf{Sorry.} This will likely be proved by showing in-degree $=1$ and
out-degree $=1$ for each factor, or by identifying an explicit inverse rule.
\end{proof}

\section{Hamiltonicity of the graph-level factors}

\begin{theorem}[single-cycle theorem for each factor]
For each $i\in\{0,1,2,3,4\}$, the permutation $F_i$ is one directed cycle of
length $m^5$.

Equivalently, the spanning $1$-factor $A_i$ is a directed Hamilton cycle.
\end{theorem}

\begin{proof}[Proof status]
\textbf{Sorry.} This is another essential missing graph-level layer. The theorem
needs a real bridge from the accepted D5 return-map / boundary package to the
claim that the lifted full-torus permutation has exact period $m^5$ and only
one orbit.
\end{proof}

\begin{remark}[what a proof would likely need]
The missing proof probably has to supply some combination of:
\begin{enumerate}[label=\textup{(\arabic*)},leftmargin=2.5em]
\item a full-torus lift from boundary exactness to a global return object;
\item an exact period computation on the lifted object;
\item a proof that no smaller orbit can survive after lifting back to the
      torus.
\end{enumerate}
At present none of those is written here.
\end{remark}

\section{Edge-disjointness and coverage}

\begin{proposition}[edge-disjointness]
The arc sets $A_0,\dots,A_4$ are pairwise disjoint.
\end{proposition}

\begin{proof}[Proof status]
\textbf{Sorry.} This should follow from the exact graph-level color rule, but
that rule is not yet written in this note.
\end{proof}

\begin{proposition}[full coverage of the torus arc set]
One has
\[
A = A_0 \cup A_1 \cup A_2 \cup A_3 \cup A_4.
\]
\end{proposition}

\begin{proof}[Proof status]
\textbf{Sorry.} This should be the graph-level completeness statement saying
that the five factors cover every outgoing arc at every vertex exactly once.
\end{proof}

\begin{corollary}[counting closure]
If each $A_i$ is a directed Hamilton cycle and the five $A_i$ are pairwise
disjoint, then they form a Hamilton decomposition of $D_5(m)$.
\end{corollary}

\begin{proof}
Each $A_i$ has exactly $m^5$ arcs. So five pairwise disjoint such cycles
contain exactly $5m^5$ arcs. But $D_5(m)$ also has exactly $5m^5$ arcs,
because it has $m^5$ vertices and out-degree $5$. Therefore the five arc sets
must cover the full torus arc set.
\end{proof}

\section{Draft final proof}

\begin{theorem}[draft graph-theoretic odd-$m$ D5 theorem]
Let $m$ be odd in the final intended D5 range.

Then $D_5(m)$ admits a Hamilton decomposition into five arc-disjoint directed
Hamilton cycles.
\end{theorem}

\begin{proof}[Proof status]
Assuming the channel reduction theorem, the global bridge descent theorem, the
single-cycle theorem for each factor, and edge-disjointness, the result follows
immediately from the counting corollary.

\textbf{Sorry.} At present those imported steps are not proved in this file, so
this is still only a proof skeleton.
\end{proof}

\section{Exact remaining work}

To turn this draft into a genuine graph-theoretic Hamilton decomposition proof,
the following pieces still have to be written.
\begin{enumerate}[label=\textup{(\arabic*)},leftmargin=2.5em]
\item Replace the phrase ``odd in the final intended D5 range'' by an exact
      quantified statement in $m$.
\item Prove the channel reduction theorem from the full odd-$m$, $d=5$ torus
      problem to the best-seed channel.
\item Restate the exact graph-level rule that turns the accepted D5 witness
      package into five global successor maps $F_i$.
\item Prove that raw global $(\beta,\delta)$ exactness is exactly what makes
      those $F_i$ globally well defined.
\item Prove that each $F_i$ is one Hamilton cycle rather than a union of
      smaller cycles.
\item Prove edge-disjointness of the five factors at the graph level.
\item Either close Inputs~A--E in the same manuscript, or cite them as formal
      earlier theorems with stable theorem-level statements.
\end{enumerate}

So this file is ready as a working manuscript scaffold, but not yet as a
finished Hamilton-decomposition proof.

\end{document}
